\hypertarget{bootstrap}{}
\hypertarget{common-date-circular-block-bootstrap-for-rank-uncertainty}{%
\chapter{Common-date circular block bootstrap for rank
uncertainty}\label{common-date-circular-block-bootstrap-for-rank-uncertainty}}

\hypertarget{why-a-rank-needs-its-own-uncertainty-analysis}{%
\section{Why a rank needs its own uncertainty
analysis}\label{why-a-rank-needs-its-own-uncertainty-analysis}}

A \emph{rank} is a discontinuous function of all 25 scores. Writing the
descending rank of portfolio $i$ in terms of its posterior score $\hat s_i$
as
\begin{equation}\label{eq:ch9-rank}
R_i = 1 + \sum_{j\ne i}\mathbf{1}\{\hat s_j > \hat s_i\},
\end{equation}
it is a step function: flat almost everywhere and jumping by an integer when
two scores cross. Its gradient is zero wherever it is defined, so the delta
method (\S\ref{prelim-asymp}) cannot attach a standard error to it, and an
arbitrarily small change in one score can move a rank by several positions.
A standard error for portfolio $i$ alone therefore cannot say whether it
stays above portfolio $j$ when the shared historical shocks are resampled.
The only honest description of rank uncertainty is the sampling distribution
of the whole ranking, obtained by resampling the entire estimator
(\S\ref{prelim-boot}).

\hypertarget{why-blocks-and-why-common-dates}{%
\section{Why blocks and why common
dates?}\label{why-blocks-and-why-common-dates}}

The bootstrap approximates the sampling distribution of a statistic by
recomputing it on data resampled from an estimate of the data-generating
process (\S\ref{prelim-boot}). Three design choices adapt the plain
independent bootstrap to this dependent panel.
\begin{itemize}
\tightlist
\item
  Monthly observations are dependent, so resampling \emph{individual}
  months would destroy serial and cross-sectional structure. The
  \emph{block} bootstrap resamples contiguous runs of length $b$ instead,
  preserving dependence within a block.
\item
  A block starting at index $s$ is taken \emph{circularly}, wrapping modulo
  $T$:
  \begin{equation}\label{eq:ch9-circular}
  \mathrm{block}(s) = \bigl((s+h-1)\bmod T\bigr)+1,\qquad h=0,\dots,b-1,
  \end{equation}
  so every month --- including those near the start and end of the record ---
  is equally likely to begin a complete block, removing the edge bias of the
  non-circular version.
\item
  The \emph{same} resampled date-index vector is applied to all 25
  portfolios and all factors (a \emph{common-date} bootstrap), so
  contemporaneous cross-portfolio shocks are preserved --- the same
  dependence the joint HAC covariance of Chapter~6 models.
\end{itemize}
Consistency of the block bootstrap requires the block length to grow with
the sample but slowly: $b\to\infty$ with $b/T\to0$ as $T\to\infty$. The
canonical run fixes $b=12$ months, one year of local dependence.

\hypertarget{full-refit-algorithm}{%
\section{Full-refit algorithm}\label{full-refit-algorithm}}

\begin{verbatim}
for b in range(B):
    idx = circular_block_indices(T, block_length=12, rng=rng)
    y_b = y[idx, :]             # same dates for all portfolios
    f_b = factors[idx, :]

    alpha_b, beta_b, resid_b = fit_all_ols(y_b, f_b)
    V_b = joint_newey_west_alpha_covariance(resid_b, f_b, lags=12)
    mu_b, tau_b = fit_empirical_bayes(alpha_b, V_b)
    post_b = posterior_mean(alpha_b, V_b, mu_b, tau_b)
    rank_b = descending_rank(post_b)
    store(post_b, rank_b)
\end{verbatim}

Do not hold the original betas, covariance or hyperparameters fixed inside
the bootstrap: every stage --- regressions, joint HAC covariance,
empirical-Bayes hyperparameters, ranking --- is refitted on each resample, so
that the reported uncertainty reflects the \emph{whole} estimator and not
just the last step. Holding any stage fixed would understate uncertainty.

Across the $B$ replicates, portfolio $i$'s rank and score uncertainty are
summarised by empirical quantiles $Q_p$ of its resampled ranks $R_i^{(b)}$
and posterior scores:
\begin{equation}\label{eq:ch9-estimand}
\text{95\% rank interval}_i = \bigl[Q_{0.025}\{R_i^{(b)}\},\,Q_{0.975}\{R_i^{(b)}\}\bigr],
\qquad
\widehat{\Pr}_i(\text{top quintile}) = \frac{1}{B}\sum_{b=1}^{B}\mathbf{1}\{R_i^{(b)}\le5\}.
\end{equation}
These are resampling distributions of the complete estimator --- not Gaussian
confidence intervals and not Bayesian credible intervals --- so they can be
asymmetric and are naturally bounded at ranks 1 and 25.

\hypertarget{a-circular-index-example}{%
\section{A circular-index example}\label{a-circular-index-example}}

Take T = 10 observations labelled 0,\ldots,9 and block length L = 4. If
random starting indices are 8, 3 and 6, the circular blocks are:

\begin{equation}
(8,9,0,1), (3,4,5,6), (6,7,8,9).
\end{equation}

Concatenate and truncate to length T, giving (8,9,0,1,3,4,5,6,6,7). For
the portfolio panel, exactly this index vector selects rows of every
portfolio and every factor. The duplicated 6 is valid bootstrap
sampling; cross-sectional alignment at each selected date is preserved.

Block length trades bias against variance. Very short blocks resemble
IID resampling and break dependence; very long blocks leave few
effectively distinct resamples. Twelve months is substantively
interpretable here, but a defensible workflow also checks nearby
lengths. With B draws, the Monte Carlo standard error of an estimated
probability \ensuremath{\hat{p}} is approximately

\begin{equation}
MCSE(\hat{p}) \approx  \sqrt{\hat{p}(1-\hat{p})/B}.
\end{equation}

At \ensuremath{\hat{p}} = 0.5 and B = 5,000 this is about 0.0071, or 0.71 percentage
points. Quantile endpoints generally need more draws than posterior
means, which is why the 1,000-versus-5,000 stability comparison matters.

\hypertarget{outputs}{%
\section{Outputs}\label{outputs}}

\begin{itemize}
\tightlist
\item
  Median rank and central 95\% rank interval.
\item
  Probability of being in the top or bottom quintile.
\item
  Bootstrap score intervals.
\item
  Stability comparison across 1,000 and 5,000 draws.
\end{itemize}

\textbf{Paper result.} With 5,000 draws, rank-interval endpoints moved
by at most one rank relative to 1,000 draws, but alpha-tail endpoints
moved by as much as 27 bps/year. SMALL HiBM is point rank 1 and reaches
rank 6 in its central rank interval.

